% !PRINTED 525
\weberpage{525}
%<<CONT
liegen, dass sich Alles leicht übersehen lässt, und doch die wichtigsten Erscheinungen der Gruppenbildung dabei zu Tage treten.

Aus vier Ziffern, 0, 1, 2, 3, lassen sich 24 Permutationen bilden, und so hoch ist also der Grad der \gesperrt{Galois'schen} Gruppe der \gesperrt{allgemeinen} Gleichung vierten Grades im Körper der rationalen Functionen der Co\"efficienten (vgl. den Schluss von §. 149). Wir stellen diese Permutationen durch ihre Cyklen in folgender Weise dar, wobei 1 die identische Permutation bedeutet:
\weq{1}{\begin{gathered}
P = 1, \quad (0,1), \quad (0,2), \quad (0,3), \quad (1,2), \quad (1,3), \quad (2,3) \\
(0,1)(2,3), \quad (0,2)(1,3), \quad (0,3)(1,2) \\
(0,1,2), \quad (0,1,3), \quad (0,2,3), \quad (1,2,3) \\
(0,2,1), \quad (0,3,1), \quad (0,3,2), \quad (1,3,2) \\
(0,1,2,3), \quad (0,1,3,2), \quad (0,2,3,1) \\
(0,2,1,3), \quad (0,3,1,2), \quad (0,3,2,1).
\end{gathered}}

Wir haben also in $P$ ausser der identischen Permutation sechs Transpositionen, acht dreigliedrige, sechs viergliedrige Cyklen und drei Permutationen von der Form $(0,1)(2,3)$, die wir \gesperrt{Transpositionspaare} nennen wollen.

In $P$ ist die alternirende Gruppe vom Grade 12 als Normaltheiler enthalten. Diese Gruppe kann keine einzelne Transposition und keine viergliedrige cyklische Permutation enthalten, und es bleiben also folgende übrig:
\weq{2}{\begin{gathered}
Q = 1, \quad (0,1)(2,3), \quad (0,2)(1,3), \quad (0,3)(1,2) \\
(0,1,2), \quad (0,1,3), \quad (0,2,3), \quad (1,2,3) \\
(0,2,1), \quad (0,3,1), \quad (0,3,2), \quad (1,3,2).
\end{gathered}}

Wir wollen nun alle vorhandenen Theiler der Gruppe $P$ aufsuchen. Dabei lassen wir die intransitiven Gruppen beiseite, die alle entweder aus drei oder aus sechs Permutationen von nur drei Elementen, oder aus den Permutationen von zwei Paaren von zwei Elementen bestehen. Wir suchen also nur die \gesperrt{transitiven Theiler} von $P$ auf. Ausserdem beschränken wir uns bei einem System conjugirter Theiler auf einen Repräsentanten, aus dem die anderen ja durch Ziffernvertauschung hergeleitet werden können (§. 154, 6.).

Eine wesentliche Vereinfachung wird durch folgende beiden Sätze herbeigeführt:

1. \gesperrt{Wenn eine transitive Permutationsgruppe von vier Ziffern einen dreigliedrigen Cyklus enthält,}
%CONT>>
